\subsection{Logistic Regression}
\label{sec:logistic_regression}

\begin{table*}[htbp]
\caption{Hyperparameter Search Space for Logistic Regression Classifier.}
\label{tab:V}
\centering
\begin{tabular}{llc}
\toprule
Component / Hyperparameter & Evaluated Values & Count \\
\midrule
Feature Scaler & None, MinMaxScaler & 2 \\
Dimensionality Reduction (PCA) & None, 0.50, 0.75 & 3 \\
Regularization Strength & 0.01, 0.1, 1, 10, 100 & 5 \\
Optimization Solver & 'lbfgs', 'saga', 'newton-cg' & 3 \\
Maximum Iterations & 500, 1000, 2000 & 3 \\
\midrule
\multicolumn{2}{r}{\textbf{Total Grid Combinations}} & \textbf{270 (1,350 fits)} \\
\bottomrule
\end{tabular}
\end{table*}

Pendekatan regresi multinomial (Softmax regression) pada LR memetakan representasi fitur laten ke dalam enam kategori demografis interseksional melalui estimasi probabilitas linier terpadu \cite{ref37}. Fungsi probabilitas Softmax menghitung peluang bersyarat sampel terhadap kategori target sebagaimana didefinisikan pada \eqref{eq:7}, di mana $\mathbf{x}$ merepresentasikan vektor fitur masukan, $y$ adalah variabel label kelas, $c$ menyatakan kategori kelas target, $\mathbf{w}_c$ dan $b_c$ masing-masing merupakan vektor bobot serta bias kelas $c$, sedangkan konstanta $K = 6$ menunjukkan jumlah total kelas interseksional \cite{ref38}:
\begin{equation}
P(y = c \mid \mathbf{x}) = \frac{e^{\mathbf{w}_c^T \mathbf{x} + b_c}}{\sum_{j=1}^K e^{\mathbf{w}_j^T \mathbf{x} + b_j}}
\label{eq:7}
\end{equation}

Proses pembelajaran model meminimalkan fungsi kerugian cross-entropy dengan menyertakan penalti regularisasi $L_2$ untuk mengendalikan kompleksitas bobot pada ruang representasi berdimensi tinggi \cite{ref39}. Pencarian hyperparameter terstruktur mengevaluasi variasi kekuatan regularisasi $C$, algoritma solver optimasi (lbfgs, saga, newton-cg), batas iterasi max\_iter, serta integrasi penskalaan dan PCA, menghasilkan 270 kombinasi evaluasi atau 1,350 total fits pada protokol 5-Fold Stratified Cross-Validation yang disajikan pada Table~\ref{tab:V}.
